% Pursuing Stacks (cotes 134-2 to 134-8): archive cote 134-7, site page 37 (PDF page 38), typescript page 479.
% typescript. Draft compared: ps6.tex:1781-1849.
% AI-transcribed by Claude Opus 5.5 (claude-opus-5-5, adaptive thinking, effort high), Message Batches API, batch ps3 (msgbatch_01XdFmuVVBAcgc7EFxdtCqsj), 2026-11-12; not yet checked by eye.
\scanpage{134-7/37}
\mspage{479}

on page 452 (taken up again on page 454 and following).   Of course in
case we do'nt assume $k=\underline{Z}$ ,  writing $\underline{Z}^{(X_*(k))}$ instead of $k^{(X_*(k))}$ as
we did in (w) p. 454 now looks kind of silly,  and the idea in this \xout{k-li-}
k-linear context to take k-valued homology of $X_*(k)$ rather than $\underline{Z}$-valued
one is evident enough !  However I was confused by the misconception that
the \uline{integral} homology of $X_*(k)$ should carry k-linear structure, as
this was what I expected too for the invariants $\pi_i(X_*(k))$ (which seems
to turn out to be correct).  This misconception was corrected a few
pages later (p.457) but still I kept dragging along the silly first
member of (w) .  Anyhow, it \inserted{now just} occurs to me that except in case $k=\underline{Z}$ , it
is definitely \uline{false} that (V') is a quasi-isomorphism, except in some
wholly trivial cases.  Indeed, let $H_i$ be the first non vanishing homolo-%
gy invariant (U) of $X_*$ (or more safely still, take i=n+1), then
we definitely expect to have a canonical isomorphism \hand{of $k$-modules}
\[ \text{(V)} \qquad \pi_i(X_*(k)) \hand{\simeq} H_i(X_*) = H_i \]
but we equally have by Hopf's theorem, as the lower $\pi_j$'s of $X_*(k)$ are
zero
\[ \pi_i(X_*(k)) = H_i(X_*(k), \underline{Z}) \quad , \]
hence \xout{(posing S = $X_*(k)$)}
\[ \text{(V')} \qquad H_i \hand{\simeq} H_i(S,\underline{Z}) \quad , \quad \text{where} \quad S = X_*(k) \quad , \]
which is not compatible with the guess that $H_i \simeq H_i(S,k)$ ($\simeq H_i \hand{\otimes}_{\underline{Z}} k$ ) .
Presumably the isomorphism (V') above is induced by the first
map in (w) above, but (except for $k=\underline{Z}$ ) we should expect in (w) to
have an isomorphism only for the lowest dimensional homology groups
which are occuring in the two first members.  Anyhow, it appears after
all that this map in (w) is the more reasonable one \inserted{compared to} \deleted{than} (V') above, as
it yi\over{\unsure{le}}{el}ds an isomorphism on homology in the key dimension n+1, whereas
(V') apparently will practically \uline{never} give an isomorphism.

What is mainly lacking still in this review of the \deleted{main} expected \inserted{main}
features of schematization of homotopy types, is description of the
k-module structure on the homotopy groups
\[ \text{(W)} \qquad \pi_i(X_*) \overset{\mathrm{def}}{=} \pi_i(X_*(k)) \quad , \]
or preferably still, a direct description of those invariants as
k-modules, working within the model category $\hand{\mathcal{M}}_n(k)$.  This, as sugges-%
ted in yesterday's notes\inserted{(p. 464)}, may be achieved \uline{by developping a theory of}
\uline{Postnikov dévissage within} $\hand{\mathcal{M}}_n(k)$ \uline{and using} (V) in order to pull oursel-%
ves by the bootstraps, defining homotopy finally in terms of homology.

\margin{37}
